\documentclass[../main.tex]{subfiles}
\begin{document}
%==========================================TIÊU ĐỀ TẬP========================================
\begin{table}[h]
  \begin{adjustwidth}{-1cm}{}
    \begin{tabular}{>{\centering\arraybackslash}p{6cm}|>{\raggedright\arraybackslash}p{12.5cm}}
      \multirow{5}{6cm}
      % Cột bên trái: ÉPISODE và số tập
      {\\[-40pt]
        \flaregothic\fontsize{20pt}{20pt}\selectfont \centering
        \color{eptype}ÉPISODE\color{black}\\
        \burbank\fontsize{72pt}{72pt}\selectfont
        \color{epnum}109\color{black}\\[6pt]
        \cafeteria\fontsize{20pt}{20pt}\selectfont\color{epnum}(9-1)\color{black}
        \\[18pt]
        % \flaregothic\fontsize{14pt}{14pt}\selectfont
        % \color{epattr}BIENVENUE, \uline{\color{Banpire} Mel} !
        \color{black}
      }
       &
      % Dòng thứ nhất: tên tập tiếng Nhật
      {\fotskip\fontsize{18pt}{18pt}\selectfont \ruby{最}{さい}\ruby{速}{そく}への\ruby{憧}{あこが}れ
      }            \\[6pt]
      % Dòng thứ hai: tên tập tiếng Anh
       & {\cafeteria\fontsize{18pt}{18pt}\selectfont The Need of Speed}                            \\[4pt]
      % Dòng thứ ba: tên tập tiếng Việt
       & {\vnmsans\fontsize{18pt}{18pt}\selectfont Xe ai nhanh hơn?}                               \\[8pt]
      % Dòng thứ tư: Nhân vật xuất hiện
       & {\caviar{\fontsize{14pt}{14pt}\selectfont Nhân vật: \emp{kuon1} \emp{ryuuhou1} \emp{suisei} \emp{marine} \emp{kanata} \emp{nene} \emp{botan}}}
      \\[8pt]
      % Dòng cuối cùng: Thời gian diễn ra của tập (thường sẽ là ngày phát hành)
       & {\continuum\fontsize{16pt}{16pt}\selectfont Ngày phát hành: 06/06/2021}                   \\[18pt]
    \end{tabular}
  \end{adjustwidth}
\end{table}
%=========================================TRUYỆN CHÍNH========================================
\color{black}

\txtbx[2]
{{\color{vol} \uline{Tóm tắt:}} Một ngày hè oi ả, Suisei tự chế tàu hỏa bằng bìa các-tông và chạy vòng quanh văn phòng, kéo theo Kanata và Ryuuhou lên trời. Marine xuất hiện với xe ba bánh, Botan với xe đẩy siêu thị ``DeLorean'', và Nene đòi lại chiếc xe của mình. Cuộc tranh cãi về ``ai nhanh hơn'' nổ ra, rồi văn phòng bị nổ tung vì một ván Old Maid đầy cay cú. Kuon tổ chức cuộc đua ba vòng với giải thưởng là 30 triệu yên tiền bồi thường, và cả ba về đích cùng lúc. Mọi người cùng nhau xây dựng lại văn phòng, và lời thề ``không bao giờ chơi Old Maid nữa'' được tuyên bố.}

Ngày đầu hạ bắt đầu với ánh nắng mặt trời chiếu rực rỡ, và Chủ Nhật đầu tiên của tháng Sáu mang lại một bầu không khí thoải mái, nhẹ nhàng cho tất cả mọi người. Mặc dù văn phòng đã bật điều hòa để giữ không khí luôn mát mẻ, những tia nắng vàng vẫn xuyên qua kẽ cửa sổ, nhắc nhở mọi người về cái nóng bức đang chờ đợi bên ngoài, mở màn cho mùa hè oi nồng nhất trong năm.

Nhưng Suisei dường như chẳng hề bận tâm đến thời tiết hay bất cứ điều gì khác, ngoại trừ niềm vui đang tràn đầy trong người. Cô đang chạy vòng quanh phòng với một chiếc tàu hỏa làm bằng bìa các-tông mặc trên người, đôi mắt lấp lánh niềm hạnh phúc ngây thơ như một đứa trẻ. Trên chiếc tàu tự chế ấy, những đường vẽ màu xanh và đen được tô tỉ mỉ, trông vừa công phu vừa đáng yêu.

``Tàu của Suisei kêu tu-tu-tu\dots''

Tiếng còi tàu giả của cô gái tóc xanh vang vọng khắp văn phòng. Tại bàn làm việc, Kuon vẫn đang tập trung vào màn hình máy tính, ngón tay gõ đều đặn trên bàn phím. Nhưng trước tiếng ồn ào vui nhộn ấy, cậu không nhịn được mà quay đầu nhìn lại một chút, rồi bật cười lớn. Ngồi cạnh cậu trên ghế xoay, Ryuuhou cũng nhìn Suisei với ánh mắt tò mò thích thú, tay còn cầm lon nước ngọt vừa mở nắp.

``Trông không khác nào quay về tuổi thơ nhỉ,'' cô em gái của cậu thì thầm, vừa cười khúc khích trước cảnh tượng trước mắt, ``giống như hồi anh em mình hồi còn nhỏ ha.''

``Chẳng khác gì hai đứa con nít đang chơi đồ hàng,'' cậu Tổng bình luận, lắc đầu nhưng nụ cười vẫn không rời khỏi môi.

Nàng Ca sĩ tiếp tục chạy vòng quanh văn phòng, vẫy tay như một nhân viên nhà ga thật sự đang mời khách: ``Bạn ơi, có muốn lên tàu cùng mình không, xình xình xịch!''

Ryuuhou quay sang nhìn Kuon, nhướng mày hỏi: ``Onii, em thấy Suisei-san vui thế. Em có thể lên tàu cùng chị ấy không? Chắc sẽ có câu chuyện thú vị để kể với các bạn trong câu lạc bộ stream mất.''

\chrint

Sở dĩ việc Ryuuhou gọi anh trai mình là ``Onii'' thay vì cách gọi trang trọng ``Onii-san'' như trước đây là do khi bước sang tuổi mười lăm, cô dần nhận ra rằng khoảng cách lễ nghi trong cách gọi đó không còn phản ánh đúng mối quan hệ gắn bó giữa hai anh em nữa - dĩ nhiên không phải theo nghĩa ``côn trùng''.

Bản thân Kuon cũng chẳng bao giờ bận tâm về cách gọi của em gái, thậm chí còn cảm thấy vui khi thấy hai anh em càng ngày càng thân thiết. Trước đây, khi còn nhỏ, cô từng gọi cậu bằng cái tên vô cùng ấu thơ ``Onii-chan'', rồi dần lớn lên và chuyển sang dùng cách gọi lịch sự hơn là ``Onii-san'', và giờ là ``Onii'' - một phong cách thân mật, hợp với tính cách tomboy vốn có của mình. Đó không phải là sự thiếu tôn trọng, mà là dấu hiệu cho thấy cô đã coi người anh của mình là một chỗ dựa thực sự, gần gũi đến mức không cần phải giữ những khoảng cách khách sáo nữa. 

Đặc biệt, sau khi chứng kiến anh trai mình ngày càng phải gánh vác nhiều trách nhiệm hơn, luôn âm thầm gánh vác mọi việc để bảo vệ cả gia đình, cô càng cảm thấy cần một cách gọi thể hiện sự gắn kết thay vì sự xa cách. 

\chrint

Đúng lúc đó, Kanata vừa bước vào văn phòng. Nhìn thấy cảnh tượng lạ lùng trước mắt, cô đứng sững người một lát, sau đó quay sang hỏi người quản lý đang ngồi xoay ghế đối diện: ``Tàu\dots\ Suisei? Cái này là chuyện gì vậy, Kuon-senpai?''

\img{9-1a}

Kuon nhún vai đáp lời cô nàng Thiên sứ: ``Anh cũng không rõ lắm. Bỗng dưng em ấy muốn làm tàu, thế là tự cắt dán bìa các-tông thành cả một con tàu hỏa thế này. Đã chạy vòng quanh phòng như vậy được khoảng 15 phút rồi đấy.''

Ryuuhou vẫy tay chào Kanata, giọng đầy hài hước: ``Kanata-san, mau lên tàu kẻo lỡ chuyến! Em cũng vừa định xin đi một vòng đây.''

Nàng Thiên sứ nhìn tiền bối với ánh mắt không thể tin nổi. Nhưng trước khi kịp nói gì, Suisei đã chạy ngay đến bên cô, đôi mắt sáng rực gọi: ``Kanata! Em đến rồi này!''

Cô nắm lấy tay Kanata và hào hứng mời: ``Em có muốn lên chuyến tàu Suisei Cáu kỉnh này không?''

Kanata giật mình ngạc nhiên: ``Đấy là chị khi bị cáu kỉnh ạ?''

Kuon bình thản tiếp lời từ phía sau: ``Có thể em không tin, nhưng đúng là thế đấy.''

Ryuuhou cũng xen vào, giọng đầy tò mò: ``Suisei-san mà cáu kỉnh thì trông thế nào nhỉ? Câu lạc bộ em có một thành viên cũng hay nổi nóng, nhưng không biết có `độ' nào bằng chị không.''

Kanata vẫn chưa hết bàng hoàng, liền hỏi thêm về con tàu hỏa kỳ lạ: ``Không phải chị thường xuyên đi diễn hát sao? Sao lại có thời gian làm cái này?''

Nhưng ngay lập tức, cô nhận lại một ánh mắt sắc lạnh từ nàng Ca sĩ. Suisei trừng mắt, giọng nói đơn điệu mà lại đầy vẻ triết lý: ``Không có niềm vui nào trên đời này là kéo dài mãi được, đúng không?''

\img{9-1b}

``Mặt chị nhìn sợ quá!!'' Kanata run rẩy, mặt tái mét, tay bất giác bám chặt lấy cánh tay Kuon như tìm điểm tựa vững chắc.

Ryuuhou nhìn Suisei rồi cười lớn: ``Suisei-san, trông chị giống ác quỷ trong truyện vậy! Nhưng mà em thích cái vibe đó đấy.''

Kuon cũng bật cười thành tiếng: ``Đây mới chính là cái `cáu kỉnh' em ấy nói đấy. Sui-chan, đừng có dọa em ấy như thế.''

Vẻ mặt nghiêm nghị của nàng ``Sao chổi'' lập tức tan biến, trở lại vui tươi như thường. Cô vỗ nhẹ vào thân tàu bìa các-tông và mời hậu bối: ``Nói đùa thôi, em cũng lên thử đi. Cùng lên tàu với Sui-chan nào!''

Kanata do dự nhìn sang cậu Tổng, ánh mắt như muốn hỏi ``Anh có chắc là mọi thứ sẽ bình thường không?''. Ryuuhou đứng dậy, vỗ vai cô động viên: ``Kanata-san, sợ gì chứ! Em cũng định lên cùng chị nữa đây. Mà Onii, anh không lên tàu cùng mọi người à? Cơ hội hiếm có lắm đấy.''

Kuon đáp lại, cười khì và xua tay: ``Cứ lên tàu đi chơi đi. Anh đang bận việc, chắc không thể tham gia được.''

Nàng Thiên sứ miễn cưỡng trèo lên con tàu hỏa, trong khi Ryuuhou thì nhảy lên đuôi tàu, tay giơ cao như một hành khách nhiệt tình: ``Em xin một vé đi vòng quanh thế giới!''

Suisei hân hoan tuyên bố trước khi xuất phát: ``Con tàu Suisei đã sẵn sàng, khởi hành thôi!''

``\textit{Vậy là chúng ta có một con tàu Suisei, với `ba mẹ con' ngồi trên đó,}'' Kuon cười thầm trong lòng, nhìn Kanata ngồi bấp bênh và Ryuuhou đang tạo dáng như một du lịch bụi thực thụ.

\img{9-1c}

``Cái quái gì mình lại dính vào thế này\dots'' Kanata nhăn mặt, không thể tin được mình lại tham gia vào trò chơi kỳ lạ này.

Ryuuhou phấn khích chỉ tay về phía trước: ``Kanata-san, mau bám chặt vào! Em nghe nói tàu này chạy nhanh lắm đấy, không giữ chắc là bị văng ra ngoài mất!''

Nhưng trước khi hai người kịp chuẩn bị, Suisei đã hô to: ``Được rồi, mọi người bám chặt vào nha. Tàu Suisei sẽ bắt đầu\dots\ \textbf{rời ga!}''

Ai cũng tưởng con tàu sẽ chỉ chạy vòng quanh văn phòng như mười lăm phút trước, thì bỗng nhiên\dots

*BÙM!*

Suisei lao đi như một tia chớp xanh, tốc độ nhanh đến mức chỉ để lại một cơn lốc xoáy trong phòng. Kuon trợn mắt nhìn cảnh tượng không thể tin được, trong đầu thầm nghĩ: ``\textit{Chắc văn phòng chưa bị thổi bay đâu nhỉ\dots}'' cậu vội ôm chặt chiếc máy tính để tránh bị gió cuốn đi.

Chỉ trong chớp mắt, cả ba người đã biến mất khỏi tầm mắt của cậu, lao qua khu Akihabara như một cơn bão. Họ chạy một vòng quanh toàn bộ thành phố, vượt qua mọi giới hạn vật lý thông thường, rồi thậm chí còn\dots

\dots lượn không chỉ một, hai mà là rất nhiều vòng quanh Trái Đất.

\begin{figure}[H]
  \centering
  \begin{subfigure}[t]{0.45\textwidth}
    \centering
    \animategraphics[autoplay,loop,width=8cm]
    {30}{../../../../Image/Book9/9-1d/9-1d.}{1}{27}
    \caption{Con tàu Suisei lao đi như một tia chớp xanh\dots}
  \end{subfigure}
  \quad
  \begin{subfigure}[t]{0.45\textwidth}
    \centering
    \animategraphics[autoplay,loop,width=8cm]
    {30}{../../../../Image/Book9/9-1d/9-1d.}{28}{57}
    \caption{\dots băng qua cả con đường của phố Akiba\dots}
  \end{subfigure}
  \quad
  \begin{subfigure}[t]{0.45\textwidth}
    \centering
    \animategraphics[autoplay,loop,width=8cm]
    {30}{../../../../Image/Book9/9-1d/9-1d.}{58}{85}
    \caption{\dots và thậm chí còn lượn vòng quanh Trái Đất!}
  \end{subfigure}
  \quad
  \begin{subfigure}[t]{0.45\textwidth}
    \centering
    \animategraphics[autoplay,loop,width=8cm]
    {30}{../../../../Image/Book9/9-1d/9-1d.}{98}{116}
    \caption{Cuối cùng, họ trở về văn phòng, với một Suisei vui thích và hai hành khách lảo đảo.}
  \end{subfigure}
\end{figure}

Chưa đầy mười giây sau, cả ba đã quay trở lại điểm xuất phát.

Suisei phanh lại gấp, vỗ tay một cái rõ to: ``Chúng ta đã về ga cuối rồi! Em thấy chuyến đi thế nào?''

Kuon trố mắt nhìn cô nàng, không giấu được sự ngạc nhiên: ``Anh thực sự bất ngờ là em không phá tung cả cái văn phòng này ra đấy\dots''

Sau khi dừng lại, Kanata lập tức lảo đảo, đầu cô quay mòng mòng, thân ngả ra sau như sắp gãy đến nơi. Ryuuhou cũng không khá hơn, cô bé cố gắng đứng vững bằng cách bám vào tường, vừa thở dốc vừa cười: ``Trời\dots\ Tàu Suisei nhanh khủng khiếp thật! Chạy còn nhanh hơn cả máy bay dân dụng!''

Thấy hai hành khách không nói được lời nào, Suisei hí hửng đoán: ``À, chắc các em vui quá nên bay lên đến chín tầng mây rồi chứ gì!''

Kuon đổ mồ hôi hột, lắc đầu nhận xét: ``Anh nghĩ là hai đứa sợ quá vãi ra quần rồi đấy. Mà cũng ngạc nhiên là cả ba người không bị cháy rụi bởi ma sát khí quyển\dots''

Đúng như cậu nói, nàng Thiên sứ ngả đầu hoàn toàn ra phía sau, mắt trống rỗng như vừa nhìn thấy được chân lý của vũ trụ sau chuyến đi kinh hoàng. Cô em gái của cậu Tổng thì đứng tựa vào tường, thở hổn hển nhưng vẫn cố nở một nụ cười mỉa mai: ``Em mà mang cái tốc độ này đi tập livestream với các bạn trong câu lạc bộ, chắc bọn em bay luôn cả phòng khi lên sóng mất.''

Kuon liếc nhìn hai người một thoáng, rồi lẩm bẩm một mình: ``Thậm chí Kana-chan còn kịp nhìn thoáng qua quê nhà của mình nữa kìa\dots''

Bất chợt, một giọng nói vang lên từ phía cửa, trộn lẫn giữa thắc mắc và chút bực bội: ``Cái trò trẻ con vô duyên gì thế này?''

Một bước chân thong thả bước vào văn phòng. Suisei đứng sững người, mắt tròn thinh trước người mới xuất hiện; một đối thủ xứng đáng để thách thức chiếc tàu hỏa của mình.

``L-Là em ấy sao!? Chiếc xe đạp của Marine!''

\img{9-1e}

Đúng như dự đoán, đó chính là nàng Thuyền trưởng Marine, đang lướt trên chiếc xe đạp ba bánh màu đỏ - chính mẫu xe cô từng dùng để đi dưới cơn mưa rào ngày đông năm ngoái\footnote{Xem lại {\flaregothic ÉPISODE 93}, quyển số Bảy.}. Ngồi phía sau bàn làm việc, Ryuuhou nhìn thấy cảnh tượng thì nhếch mép, giọng đầy trêu chọc: ``Chiếc xe ba bánh màu đỏ này thực sự hợp với Marine-san. Nhưng mà sao chị lại đi xe đạp thay vì dùng tàu thuyền như thường lệ?''

Kuon nhướng mày nhận xét: ``Bản thân em cũng chơi đồ trẻ con thế mà còn dám chê người khác\dots\ Tưởng đi được xe cộ hạng sang gì chứ.''

Kanata cúi nhìn phương tiện của Thuyền trưởng, mặt nhăn lại: ``Đó là xe ba bánh mà!''

Cậu Tổng chống cằm, mỉm cười khà khà: ``Chắc là em ấy bắt chước Luna-chan quá nhiều rồi đấy\dots''

Em gái của cậu cũng gật gù, giọng vẫn mang chút mỉa mai: ``Nếu Luna-san thấy Marine-san đạp xe này, chắc chị ấy sẽ bị chọc cười đến thối mũi mất.''

Marine lướt nhanh qua bên cạnh chiếc tàu bìa các-tông của Suisei, ánh mắt lấp lánh niềm tự hào. Cô dừng xe trước mặt Kanata, nháy mắt hỏi: ``Kanatan, em thấy xe của chị có ngầu không?'' Thậm chí, đôi mắt cô còn sáng rực, đầu lắc lư không ngừng như một đứa trẻ đang mong chờ được khen ngợi.

Nàng Thiên sứ chớp mắt vài lần, liếc qua chiếc xe ba bánh đỏ rực của đàn chị. Sau vài giây quan sát, cô nghiêng đầu, trả lời thẳng thắn: ``Nhưng đó là \uline{xe ba bánh} mà.''

Không khí trong phòng lập tức trở nên căng thẳng.

Ryuuhou khoanh tay, khẽ cười: ``Kanata-san nói thật quá đấy. Nhưng mà cũng đúng thôi, xe ba bánh thì vẫn là xe ba bánh.''

Mắt Marine trợn tròn, gân xanh nổi lên trên trán. Cô lập tức lao tới, ghé sát mặt vào đàn em, giọng gắt gỏng như một thuyền trưởng thực thụ đang quát mắng thủy thủ: ``Chị biết rồi mà, cái thiên sứ cứng đầu này! Lần đầu tập đạp xe thì không dùng xe ba bánh thì dùng cái gì hả?''

\img{9-1f}

Kanata giật bắn mình, toàn thân run lẩy bẩy: ``Eek! Em xin lỗi!''

Kuon nhìn cảnh tượng trước mắt, khoanh tay lắc đầu ngao ngán: ``Giờ anh còn chẳng biết ai mới là người đáng sợ hơn nữa\dots''

Ryuuhou nhìn Marine với vẻ thán phục giả tạo: ``Marine-san gắt gỏng y hệt một hải tặc thực thụ luôn. Chỉ tiếc là thứ chị cầm lái không phải tàu chiến mà lại là\dots\ xe con nít.''

Senchou vẫn chưa chịu dừng cuộc tranh cãi, tiếp tục chỉ trích món đồ chơi của Suisei: ``Với cả, cái tàu của em cũng chỉ là mấy tấm bìa các-tông ghép lại thôi!''

Nàng Ca sĩ khoanh tay, nhếch môi đầy tự hào: ``Được thôi, nhưng chị đây còn có `bàn tay vàng' Midas\footnote{Một trong ba thành viên của hoàng gia Phrygia trong thần thoại Hy Lạp. Ông được nhớ đến vì khả năng có thể chạm mọi thứ thành vàng. Suisei sử dụng thần thoại đó để lòe đối thủ của cô.} nữa kìa.''

Cả Kuon và Kanata cùng nghiêng đầu, đồng thanh hỏi: ``Cái gì cơ?''

Ryuuhou cũng nhíu mày, giọng đầy nghi ngờ: ``Suisei-san, Midas liên quan gì đến chiếc tàu của chị vậy ạ? Hay chị định biến cái tàu bìa này thành vàng thật?''

Marine khoanh tay, cười khẩy: ``Hay đấy. Nhưng em đây còn có siêu tốc độ Flash nữa cơ!''

``Hả?''

Suisei hất mái tóc xanh ngược ra sau, ánh mắt sắc bén đầy thách thức: ``Cô em nghĩ mình nhanh hơn chắc? Chị đây vừa chạy vòng quanh Trái Đất 20 vòng chỉ trong 10 giây thôi đấy!''

Đối thủ của cô đập mạnh tay lên tay lái xe ba bánh, cười hả hê: ``Chị nghĩ thế là giỏi á? Em đây còn đi vòng quanh Dải Ngân Hà 10 vòng trong cùng 10 giây á!''

Nhìn hai cô nàng đấu khẩu nhau để chứng minh đồ chơi của mình đỉnh hơn, chẳng khác nào hai đứa trẻ mẫu giáo, Kuon xoa thái dương, thở dài chán nản: ``\hl{Hai đứa dở hơi này lại cãi nhau vì mấy chuyện trẻ con vớ vẩn này à?}'' Cô em gái Ryuuhou cũng cũng lắc đầu, nét mặt đầy bất lực không kém gì ông anh: ``Marine-san, Suisei-san, hai chị đang làm gì thế này? Cãi nhau xem ai có món đồ chơi ngầu hơn ạ?''

Kanata cũng khoanh tay, lắc đầu ngao ngán: ``Thật sự. Nhảm nhí không thể tưởng tượng nổi.''

Cô nàng Trưởng nhóm quay sang nhìn Kuon, nhếch mép cười: ``Onii, thấy hai chị ấy làm em nhớ hồi nhỏ, hai anh em mình cũng hay cãi nhau xem ai có nhiều đồ chơi hơn đấy.''

Kuon đưa tay xoa đầu em gái, cười khổ: ``Y chang luôn.''

Khi hai cô nàng vẫn đang cãi nhau không ngớt về tốc độ và những khả năng ``siêu phàm'' mà họ tự cho mình có, bỗng một giọng nói chắc nịch vang lên từ cửa phòng, mang theo một khí thế khó tả: ``Đủ rồi, đừng tranh cãi nữa.''

Tiếng bánh xe kẽo kẹt lăn dần vào văn phòng khiến cả bốn người đang có mặt đều đồng loạt quay đầu. Ai cũng há hốc mồm kinh ngạc trước cảnh tượng trước mắt, đặc biệt là Suisei và Marine - cả hai lập tức giật mình, lắp bắp cùng một lúc: ``L-Là cô ấy! Tay đua Botan!''

Một đối thủ xứng đáng khác chính thức gia nhập cuộc chơi. Botan, cô nàng Sư tử, đang ngồi gấp chân trên chiếc xe đẩy hàng siêu thị - cái đồ mà cô vừa ``mượn'' được từ một trò chơi khiến cô cười đến chảy cả nước mắt\footnote{Tác giả đang nói đến trò chơi \textit{Trials Rising}.}.

\img{9-1g}

Ryuuhou chống cằm chăm chú nhìn chiếc xe đẩy kỳ lạ, rồi nhướng mày hỏi: ``Botan-san, chị lấy cái xe này ở đâu vậy ạ? Em tưởng mấy siêu thị không cho phép khách mang xe đẩy đi ra ngoài đâu.'' Giọng nói của cô bé pha lẫn vừa thích thú vừa nghi ngờ.

Hai người còn lại trong phòng cũng không thể che giấu sự bất ngờ và kinh ngạc của mình.

Kuon nheo mắt nhìn ``phương tiện đua'' của cô nàng tay đua Sư tử, cứng họng mấy giây mới nói nên lời: ``Vãi cả\dots\ Xe đẩy hàng siêu thị á!?''

Kanata cũng chẳng kém cạnh gì, mở to mắt hỏi: ``Chị kiếm cái này ở đâu ra thế!?''

Botan khoanh tay, liếc qua hai ``chiến mã'' của hai cô nàng tự xưng là tay đua kia, rồi nhếch mép cười khinh khỉnh: ``Xe của hai chị trông bèo nhèo thế, làm sao mà đua được. Nếu không biết cách độ xe cho đúng, thì đừng mơ tới việc đánh bại tôi.''

Ngay lập tức, Kuon đổ mồ hôi hột, tay xoa thái dương đầy bất lực: ``Độ cái gì cơ?''

Ryuuhou nghiêng đầu, nhìn cô nàng Sư tử với vẻ tò mò xen lẫn một chút chọc ghẹo: ``Xe đẩy siêu thị mà chị cũng độ được nữa hả? Chị định lắp động cơ vào nó à?''

Kanata cũng bối rối không kém, hỏi nhỏ: ``Thắng cái gì cơ? Chị cũng không hiểu gì luôn à.''

Nghe xong lời tuyên bố gây sốc từ đàn em, Suisei và Marine lập tức ôm chầm lấy phương tiện của mình, than vãn khóc lóc như thể vừa bị ai đó xúc phạm nặng nề.

Nàng Ca sĩ ôm chặt con tàu bìa các-tông tự làm của mình, giọng đầy oan ức: ``Nhưng chị vất vả bao lâu mới vẽ được chiếc tàu này mà!''

Nàng Thuyền trưởng thì ôm lấy chiếc xe ba bánh, gục đầu lên yên xe như vừa bị ai đó cười nhạo: ``Chiếc xe ba bánh này của chị đắt lắm đấy chứ!''

\img{9-1h}

Nhìn hai cô nàng đang khóc lóc ôm lấy đồ chơi của mình, Ryuuhou vừa thấy thương hại lại vừa buồn cười. ``Hai chị trông y hệt mấy đứa trẻ bị người ta cướp mất kẹo ấy. Botan-san đâu có nói gì quá đáng đâu mà,'' cô bé vừa nói vừa che miệng cười, mắt híp lại vì vui.

Kuon cũng phì cười nhìn cảnh tượng dở khóc dở cười trước mắt, thở dài: ``Thật là hai đứa con nít. Mới bị chê một câu là đã dỗi hờn.''

Kanata thì vẫn còn đứng hình, nhìn hai người đang ôm chặt phương tiện của mình mà thắc mắc: ``Sao họ trông như bị xúc phạm ghê thế!?''

Ryuuhou quay sang nàng Thiên sứ, giọng đùa giỡn: ``Chị ơi, bởi hai chị ấy dồn cả tâm huyết vào mấy cái xe đồ chơi đó chứ sao. Còn như em, em chỉ dồn cảm xúc vào mấy cây đàn của mình thôi.'' Vừa nói xong, em liếc nhanh qua Kuon rồi lập tức quay mặt đi, như thể vừa lỡ nói ra một bí mật nào đó.

Ngay khi màn than vãn vẫn chưa kết thúc, lại thêm một người nữa bước vào văn phòng. Nhưng lần này, người đến không phải để tham gia cuộc đua\dots\ mà là đến đòi lại đồ của mình!

Nene hét to bằng giọng đầy hùng hổ: ``\textbf{Tôi bắt được bà rồi, Shishiron!}''

Botan giật bắn mình, lập tức co người lại y hệt đứa trẻ bị bắt quả tang trộm bánh ăn. Cô quay đầu, mặt tái mét khi nhận ra người vừa xuất hiện: ``Gah, Momosuzu Nene!?''

Nàng Hải cẩu chỉ thẳng vào cô bạn Sư tử, quát lớn: ``\textbf{Trả lại DeLorean của tôi ngay!}'' 

Botan lắc đầu quầy quậy như đứa trẻ bướng bỉnh: ``Không trả!''

Cô nàng Trưởng nhóm khoanh tay đứng nhìn, vẻ mặt thích thú trước cảnh tượng vừa diễn ra: ``Nene-san, DeLorean mà chị nói là cái xe đẩy siêu thị này á? Trông cũng ngầu đấy, nhưng em thấy nó có giống gì so với với xe thời gian đâu.''

Kuon cũng nhíu mày chăm chú nhìn chiếc xe đẩy mà Botan đang đứng cạnh. Cậu chỉ vào nó, không thể tin được: ``Ặc\dots\ Cái xe đẩy đó mà là DeLorean á?''

Nene gật đầu lia lịa, mặt vẫn hằm hằm: ``Chứ còn gì nữa! Của em mà!'' Rồi cô quay lại trách mắng Botan: ``Này Shishiron, bà có trả cho tôi không thì nói!''

Nàng Sư tử vẫn kiên quyết lắc đầu: ``Không!''

Nàng Hải cẩu chống nạnh, bặm môi: ``Tôi không quan tâm! Tôi cần cái xe đó để đi trúng xổ số bây giờ!''

Nghe vậy, Ryuuhou nhíu mày hỏi: ``Trúng xổ số bằng xe đẩy á? Nene-san, em không hiểu lắm, chị định làm gì với nó thế?'' Cô bé nghiêng đầu, vừa hoang mang vẫn còn giữ được nụ cười tinh nghịch.

Càng nghe càng thấy mọi thứ càng trở nên vô lý, nhưng vì đã quá quen với những màn hài hước của mọi người, Kuon chỉ còn cách thở dài chán nản: ``Hai người này đang diễn vở kịch gì thế không biết\dots''

Botan vẫn ôm chặt lấy tay lái xe đẩy, hét lên như đứa trẻ bị người ta giành mất đồ chơi: ``\textbf{Không chịu đâu!!!}''

Ryuuhou nhìn Botan, giọng giả vờ thương hại: ``Botan-san, chị cứ thế này thì Nene-san sẽ thật sự đẩy xe đi mất đấy.'' Mặc dù nói vậy, trong mắt cô bé lại ánh lên niềm vui khó giấu trước cảnh tượng dở khóc dở cười.

Dù Botan có phản đối ra sao, Nene vẫn lẳng lặng đẩy chiếc xe đẩy mang tên ``DeLorean'' ra khỏi văn phòng, bỏ lại sau lưng mọi người đang đứng ngẩn ngơ, không thể tin được những gì vừa xảy ra cùng với tiếng la hét thất thanh của một đứa trẻ đòi quà Botan.

\img{9-1i}

Chứng kiến toàn bộ câu chuyện từ đầu đến cuối, Kanata vẫn chưa hết bối rối. Cô chớp mắt vài lần rồi buột miệng hỏi: ``Chuyện gì đang xảy ra ở đây thế nhỉ?''

Kuon nhún vai, ngán ngẩm đáp: ``Biết chết liền à.''

Ryuuhou đứng dậy, phủi nhẹ bụi trên quần, rồi nói bằng giọng hài hước: ``Onii này, em nghĩ hôm nay các chị ấy đang tổ chức một buổi triển lãm xe độc đáo quá. Có tàu hỏa, xe ba bánh, xe đẩy siêu thị - thiếu mỗi xe tăng là đủ cho một cuộc diễu binh hoàn chỉnh rồi.''

Sau khi Nene đẩy chiếc xe hàng, với Botan vẫn đang nhõng nhẽo đeo bám, biến mất khỏi tầm mắt, Suisei và Marine trao nhau một cái nhìn, ánh mắt bỗng lóe lên tia hiếu chiến không thể che giấu. Dù vừa tiễn đi một đối thủ đáng gờm, mối thù lâu đời về ai mới là người sở hữu tốc độ tuyệt đỉnh giữa họ vẫn chưa hề phai nhạt.

Suisei khoanh tay trước ngực, nghiêng đầu một chút như đang suy tư điều gì, rồi cất tiếng: ``Chúng ta đâu thể cứ đứng đây lườm nhau mãi được, đúng không?''

Nghe vậy, Marine nhếch mép môi, ngước đôi mắt đầy vẻ thách thức lên nhìn đối thủ: ``Sao vậy? Chị định giải quyết sao đây?''

Suisei chỉ thẳng ra phía cửa lớn của văn phòng, mạnh dạn đưa ra đề xuất: ``Chúng ta ra ngoài kia, và giải quyết mâu thuẫn này bằng một cuộc đua!''

Marine lập tức đập mạnh tay lên yên xe ba bánh của mình, gật đầu một cách dứt khoát: ``Đồng ý! Chơi luôn!''

Cả hai người tiếp tục trao nhau những ánh mắt lườm nguýt, toàn thân như thể đang bốc lên những ngọn lửa của sự cạnh tranh gay gắt, trông chẳng khác nào sắp ăn tươi nuốt sống đối phương đến nơi. Họ còn không ngừng phụ họa bằng những động tác vung tay nhún chân, y hệt như hai đứa trẻ đang khè nhau trước khi bắt đầu một cuộc đua.

Nhìn cảnh tượng ấy, Ryuuhou khoanh tay đứng một bên bật cười khẩy: ``Hai chị đang chuẩn bị đấu võ mồm à? Em thấy không khí ở đây còn sôi động hơn cả buổi tập nhạc chung của câu lạc bộ nữa.''

Còn Kuon thì đứng cạnh em gái, một tay chống cằm, mắt khẽ nheo lại với vẻ mặt thản nhiên của một người đã quá quen thuộc với những trò ngớ ngẩn không hồi kết của các đồng nghiệp. Cậu thở dài một tiếng nặng nề, rồi mới cất giọng: ``Anh cũng rất muốn tham gia vào cuộc đua của các em, mà thật không may là\dots''

Cậu khẽ hất đầu về phía cửa sổ, nơi mưa vẫn đang đổ xuống không ngừng. Kanata cũng lập tức chỉ tay về hướng đó, giọng có chút ngập ngừng như đang cố gắng nhắc nhở hai người: ``Ừm\dots\ ngoài kia trời đang mưa to lắm, hai chị ơi.''

Ryuuhou lúc này đang đứng bên cửa sổ, mắt nhìn ra ngoài vào màn mưa xối xả đang trút xuống thành phố, rồi quay mặt lại nhìn hai người đàn chị: ``Suisei-san, Marine-san, các chị định vẫn muốn ra ngoài đua xe trong tiết trời thế này à?''

Cơn mưa đến bất chợt này, thật trớ trêu thay, đã lập tức dập tắt hoàn toàn bầu không khí hừng hực chiến ý đang bùng cháy trong hai kẻ hiếu thắng kia.

Suisei bĩu môi dài ra, giọng đầy hậm hực: ``Trời ạ, thế thì chúng ta không thể đua được nữa rồi.''

Marine cũng chẳng khá hơn, cô thở dài ngao ngán, gục cằm lên tay lái chiếc xe ba bánh của mình: ``Chán thật sự, đúng là hết hơi.''

Nhưng chỉ vài giây sau, Suisei bỗng nảy ra một ý tưởng mới, mặt lập tức tươi tỉnh trở lại và hào hứng đề xuất: ``Thôi được, nếu không đua được thì chúng ta chơi Tiềm Ô Quân (Old Maid) vậy!''

Marine suy nghĩ ngẫm lại trong vài giây, rồi cũng gật đầu đồng tình với đề xuất này: ``Nghe cũng ổn đấy, được thôi!''

Cả hai cùng quay sang nhìn Kuon, trong mắt lấp lánh ánh mắt mong chờ được đồng hành: ``Kuon-senpai, chơi cùng bọn em đi nào!''

Ryuuhou cũng hùa theo, quay sang nhìn anh trai mình với giọng đầy háo hức không thể kìm nén: ``Onii, chơi cùng mọi người đi! Em cũng muốn thử sức với hai chị Suisei và Marine lắm!''

Thế nhưng\dots\ người mà tất cả mọi người đang mời gọi đã biến mất từ lúc nào không hay.

Kanata chớp mắt vài lần, ngó quanh quất khắp căn phòng rồi thắc mắc với sự ngạc nhiên: ``Ủa? Kuon-senpai đâu hết rồi?''

Marine cũng gãi đầu khó hiểu, phỏng đoán: ``Hay là anh ấy có việc bận đột xuất nên đi trước nhỉ?''

Suisei chỉ nhún vai tỏ vẻ không bận tâm: ``Cũng không quan trọng nữa. Thôi thì chúng ta bắt đầu thôi!''

Ryuuhou lập tức ngồi xuống, giơ tay lên như một đứa trẻ đang xin phép tham gia trò chơi: ``Em xin một suất! Em thay Onii em chơi vậy, được không?''

Ngay sau đó, Suisei liền lấy ra một xấp bài được xếp gọn gàng từ trong túi áo, trên mặt nở một nụ cười roi rói đầy tinh nghịch. Cả nhóm quyết định để Kanata làm trọng tài cho ván bài. Ryuuhou thì hào hứng xung phong nhận nhiệm vụ trộn bài, đôi mắt sáng rực lên đầy thích thú: ``Để em trộn bài cho mọi người nhé! Hồi nhỏ Onii hay chơi bài với em lắm, em trộn bài siêu khéo!''

\ct

Cậu con trai trưởng của gia đình Hikawa, Kuon, lúc này đang chạy bán sống bán chết dưới màn mưa tầm tã, trông chẳng khác nào vừa trốn thoát khỏi một cơn ác mộng khủng khiếp.

``\textit{Mình đã từng dính vào một ván bài thảm hại với Noel-chan, Suba-chan, Lamy-chan và Wata-chan rồi\dots\footnote{Xem lại {\flaregothic ÉPISODE 99}, quyển số Tám.} Không đời nào mình lại dẫm vào vết xe đổ lần nữa! Thôi thì ba mươi sáu kế, chuồn là thượng sách!}''

Dưới cơn mưa vẫn đang xối xả, cậu cố gắng chạy thật xa khỏi văn phòng nhất có thể, sợ rằng nếu chậm chân một chút, cậu sẽ lại bị cuốn vào bất kỳ trò chơi ``ngu xuẩn'' nào nữa của những người đồng nghiệp chẳng bao giờ lớn được này.

Khi vừa chạy được một quãng đường nhất định, cậu bỗng bắt gặp một cảnh tượng khó tin trước mắt: Nene đang vui vẻ đẩy chiếc xe đẩy hàng siêu thị đi dưới mưa, và ngồi chễm chệ trên chiếc xe đó không ai khác chính là Botan, người đang trông ỉu xìu hẳn đi vì bị cô bạn kiểm soát chặt chẽ.

Kuon chớp mắt vài lần liên tiếp, tự hỏi liệu mình có đang bị mưa làm cho hoa mắt hay không. Nhưng trước khi kịp suy nghĩ thêm bất cứ điều gì, một âm thanh kinh hoàng đột ngột vang lên từ phía sau lưng cậu.

  {\Huge \textbf{*ẦM!!!*}}

Cả mặt đất như rung chuyển theo vụ nổ đó. Một cột lửa khổng lồ lập tức bùng lên từ hướng văn phòng \hll.

Kuon dừng bước ngay lập tức, đưa tay lên che mắt để nhìn về phía cột lửa khổng lồ đang bốc cháy, trong khi Nene và Botan thì há hốc mồm vì sốc trước cảnh tượng vừa xảy ra.

Nàng Hải cẩu chớp mắt liên tục, lắp bắp hỏi với vẻ không thể tin được: ``C-Chuyện gì vừa xảy ra vậy, Kuon-senpai!?''

Cậu Tổng thở dài, gương mặt chẳng có chút ngạc nhiên nào, như thể đã quá quen thuộc với những tình huống bất ngờ kiểu này từ lâu: ``\dots Lại nữa rồi.''

Nene quay sang nhìn cậu với ánh mắt đầy ngạc nhiên, hỏi lại: ``Ý anh nói `lại nữa rồi' là sao ạ?''

``Văn phòng lại nổ nữa rồi. Kiểu này lại thêm một núi tiền phải dùng để sửa chữa nữa đây.''

Botan lúc này cũng tròn mắt hết cỡ, nhìn chằm chằm vào Kuon với vẻ khó tin: ``Đây không phải là lần đầu tiên cả văn phòng bị nổ sao anh?''

``Cũng được lần thứ năm, thứ sáu gì đó rồi,'' cậu trả lời, trên mặt nở ra một nụ cười méo xệch đầy bất lực. Cả Nene và Botan đều rơi vào trạng thái sốc tận óc trước câu trả lời bình thường đến lạ của cậu, đến mức không ai nói nên lời với nhau.

Nhưng điều kỳ lạ hơn cả là\dots\ chỉ vài phút sau vụ nổ kinh hoàng đó, trời bỗng dưng quang đãng trở lại một cách thần kỳ. Cơn mưa lớn đã dừng hẳn, ánh mặt trời dần ló rạng phía sau những đám mây đen đang tan biến.

Nene nheo mắt nhìn lên bầu trời trong xanh vừa xuất hiện, không giấu được sự ngạc nhiên: ``Hửm? Trời sao tự dưng tạnh mưa thế này?''

Kuon cũng ngước nhìn lên cao, và bắt gặp trên bầu trời cao vút, bốn bóng người quen thuộc đang lơ lửng giữa những tia nắng vàng rực rỡ. Đó là Marine, Suisei, Kanata và\dots\ Ryuuhou, cả bốn người đang cười nói vui vẻ, tay vẫn còn cầm mấy lá bài chơi Tiềm Ô Quân lúc nãy.

\img{9-1j}

Ryuuhou từ trên cao vẫy tay xuống đất, gọi lớn về phía cậu: ``Onii! Bọn em vừa thắng một ván bài cực hay luôn! Mà\dots\ hình như văn phòng bị nổ rồi hả anh?''

Kuon đứng lặng người nhìn cảnh tượng khó tin trước mắt, rồi thở dài lẩm bẩm một mình: ``Mùa mưa năm nay kết thúc sớm nhỉ\dots''

Nene quay sang nhìn cậu, chớp mắt đầy khó hiểu trước câu nói ấy: ``Nhưng mới chỉ là đầu tháng Sáu thôi mà?''

Dù sao đi nữa, khuôn mặt rạng rỡ của bốn người trên trời cho thấy họ thực sự đã làm lành với nhau, dường như chẳng còn ai để tâm gì đến cuộc chiến tranh tốc độ vô nghĩa lúc đầu nữa\dots

%==========================================TRUYỆN PHỤ=========================================
\omk

\dots hoặc ít nhất, là Kuon nghĩ vậy.

Ngay sau vụ văn phòng \hll{} lại một lần nữa hóa thành khói lửa, hai kẻ thích gây chuyện Suisei và Marine lập tức lao vào một cuộc cãi lộn không hồi kết, tranh luận mãi không dứt xem ai mới là người phải gánh trách nhiệm đền bù thiệt hại.

Nàng Thuyền trưởng khoanh tay trước ngực, mặt đầy vẻ mừng thầm, chỉ thẳng vào mặt Suisei mà đổ tội: ``Hừ! Chẳng phải là cô thua ván bài Old Maid nên cay cú, tức giận đến mức đốt cả công ty đi sao? Nhìn mặt cô là tôi biết hết rồi đấy!''

Nàng Ca sĩ trợn tròn mắt, giọng nói bỗng dưng cao vút lên: ``Cái gì vậy!? Rõ ràng là cô trước đó cứ chọc tức tôi, làm tôi mất tập trung nên mới thua! Mọi lỗi đều thuộc về cô!''

Marine chống tay vào hông, nhướng mày cười nhạo: ``Ối giời\~{} Ai bảo mà tâm lý yếu đuối thế chứ?''

Suisei siết chặt nắm tay, gân xanh bắt đầu nổi dần lên trên trán: ``\textbf{Cô vừa nói cái gì hả bà già khó tính này!?}''

``\textbf{Cô dám gọi tôi là gì cơ!?}''

Hai người dần dần tiến lại gần nhau hơn, ánh mắt như có thể bắn ra lửa. Không khí trong phòng trở nên căng thẳng đến mức cả Botan và Nene đứng xem cũng cảm nhận được những tia sét vô hình đang chớp giật giữa hai người.

Ở ngoài cuộc chiến lời qua tiếng lại, Kuon chỉ có thể đứng nhìn hai cô gái cãi nhau mà lắc đầu ngao ngán: ``Thật sự là hai đứa con nít, không lớn được bao giờ.''

Kanata đứng bên cạnh cậu, cũng gật gù đồng tình: ``Đúng thật\dots\ Trông y hệt hai đứa trẻ vừa bị cướp mất kẹo.''

Ryuuhou, người vẫn còn thở hổn hển sau chuyến bay ngắm cảnh vừa rồi, bước lại gần anh trai mình, giọng đầy hài hước: ``Onii, em cứ nghĩ các chị ấy đã nguôi giận sau khi được lên trời chơi đùa. Ai ngờ vừa đặt chân xuống đất là lại tiếp tục cãi nhau rồi.''

Hai người cùng đứng ngoài quan sát, chẳng có ý định can thiệp vào cuộc cãi lộn của họ. Mãi cho đến khi thấy hai kẻ hiếu thắng kia bắt đầu xắn tay áo, chuẩn bị lao vào ``đấu tay đôi'' thật sự, thì hai người em của Thế hệ thứ Năm mới vội vàng xông tới kéo họ ra, tránh cho mọi chuyện càng đi càng xa.

Nhưng dù đã bị ngăn cản, cả hai vẫn không chịu dừng lời qua tiếng lại, chẳng ai chịu nhường ai, y hệt hai con dê đang tranh nhau đi qua cây cầu hẹp, chẳng ai chịu lùi bước.

Cuối cùng, không chịu được cảnh tượng ấy nữa, Kuon đứng dậy, vỗ tay cái *BỐP!* thật to để thu hút sự chú ý của tất cả mọi người: ``Được rồi, dừng lại hết! Không ai chịu đền tiền sửa cái công ty này đúng không?''

Hai người đang cãi nhau bỗng dưng im lặng lập tức, quay sang nhìn cậu với ánh mắt chăm chú. Cả hai đồng loạt gật đầu cái rụp, đồng thời vẫn không quên trừng mắt nhìn nhau thêm cái.

Thấy tình hình đã nằm trong tầm kiểm soát, cậu Tổng hít một hơi thật sâu, nở ra một nụ cười đầy vẻ ẩn ý rồi lớn tiếng tuyên bố: ``Vậy thì chúng ta sẽ tổ chức một cuộc thi để phân định ai thắng ai thua!''

Tất cả mọi người đồng loạt quay sang nhìn cậu với ánh mắt đầy nghi ngờ.

Nene chớp chớp mắt, hỏi nhỏ: ``Cuộc thi gì thế ạ?''

``Một cuộc đua \textbf{Ai nhanh hơn!}''

Tất cả năm cô gái trong phòng đều sững sờ trước lời thách đấu bất ngờ ấy. Ngay cả Suisei và Marine, hai người vừa mới gân cổ cãi nhau nảy lửa trước đó, cũng không ngờ Kuon lại đề xuất ra chuyện đó.

Bỗng một tiếng hét ngạc nhiên từ trên cao vọng xuống đất: ``\textbf{Ể!?}''

Ryuuhou đứng cạnh Kanata, giơ tay lên như đứa trẻ xin được tham gia trò chơi: ``Onii, em có được tham gia không ạ? Em cũng muốn thử sức với các chị ấy!''

Kuon nhìn em gái mình, nhếch mép đáp: ``Khỏi. Để mấy chị ấy tự giải quyết chuyện của họ đi.''

Ryuuhou bĩu môi tỏ vẻ thất vọng, nhưng rồi cũng cười nhượng bộ: ``Thôi được, vậy em làm trọng tài phụ nhé.''

\ct

Trên một con phố hoang vắng không có bóng dáng một ai, ba ``tay đua bất đắc dĩ'' đã sẵn sàng vào vị trí của mình. Đây là vùng ngoại ô ít người qua lại, với những cung đường rộng lớn thoáng đãng, hoàn toàn lý tưởng cho một cuộc đua mang đầy tính chất\dots\ tàn phá. Vạch xuất phát được dựng ngay trước một siêu thị lớn, chỉ cách văn phòng không xa.

Ryuuhou đứng bên cạnh Kuon, tay cầm một tờ giấy nháp đang ghi chú nhanh vài con số, đôi mắt tinh nghịch lướt qua từng chiếc xe xếp hàng. Cô vừa liếc nhìn chiếc xe ba bánh rung lắc của Marine, vừa khẽ viết bậy vài dòng không mấy tốt đẹp lên góc giấy, miệng nén cười để không lộ vẻ chọc ghẹo.

Ba chiếc ``chiến mã'' cũng đã sẵn sàng xếp hàng:

\begin{itemize}
  \item Suisei ngồi trên chiếc ``Tàu hỏa'' độc đáo của mình: một thùng bìa các-tông đeo quanh người, được cô trang bị thêm vài chiếc tên lửa nhỏ để tăng ``hiệu ứng'' .
  \item Marine tự tin chễm chệ trên chiếc xe đạp ba bánh cỡ người lớn, trông tưởng chắc chắn mà thực chất lại rung lắc bất ổn.
  \item Còn Botan thì đứng thẳng trên một chiếc xe đẩy hàng siêu thị, với Nene đóng vai trò là ``người đẩy'' tạm thời cho cô.
\end{itemize}

Không ai trong số họ là những tay đua chuyên nghiệp, thế nhưng ai cũng ấp ủ trong mình ngọn lửa khát khao chiến thắng\dots\ hay ít nhất là trốn tránh được nghĩa vụ bồi thường hàng chục triệu yên.

Ryuuhou bước ra giữa vạch xuất phát, giơ tay ra hiệu cho mọi người yên lặng: ``Bây giờ Onii - Hikawa Kuon - sẽ phổ biến luật đua cho tất cả mọi người. Mong các chị lắng nghe ạ.''

Tất cả lập tức dồn toàn bộ sự chú ý vào cậu Tổng. Cô cũng đóng tậo giấy lại, đứng bên cạnh Kanata với tư thế trọng tài chuyên nghiệp, chỉ còn chút ánh mắt nháy nhẹ với nàng Thiên sứ để lộ sự hào hứng.

Kuon hắng giọng một lúc lâu, rồi đi về phía giữa vạch xuất phát, nơi mà cô em gái của cậu vừa mới lon ton chạy về: ``Tuyến đường đua của chúng ta sẽ bắt đầu từ đây, vòng quanh khu phố này, rẽ vào đường Hamari, chạy qua Đại lộ Mirete, rồi quay lại điểm xuất phát. Chúng ta sẽ đua tổng cộng ba vòng.''

Marine nhíu mày, giơ tay ngắt lời cậu: ``Khoan đã, vậy tổng cộng mỗi vòng chúng ta phải chạy bao nhiêu cây số?''

Cậu trả lời dứt khoát: ``135 cây số cho mỗi vòng.''

Suisei chỉ nhún vai, gật gù như không mấy bận tâm: ``Nghe có vẻ cũng không xa lắm nhỉ. Em còn chạy được cả nghìn cây số nữa là.''

Kuon chỉ mỉm cười bí ẩn: ``Hy vọng là vậy.'' Cậu tiếp lời: ``Hai em sẽ đua đối kháng với nhau, cộng sự tham gia thêm là Botan-chan--''

``Đây là xe của Nene!'' cô nàng Hải cẩu lập tức nhảy vào sửa lỗi.

``Cũng thế thôi, có khác gì đâu,'' cậu đáp đầy hời hợt. Marine lập tức phản ứng vội vàng: ``K-khoan đã! Tại sao Shishiron cũng được tham gia cuộc đua này?''

Cậu Tổng khoanh tay, trả lời một cách thản nhiên: ``Để lấy mốc thời gian thôi. Mà thêm vào là em ấy cũng muốn đua nữa. Anh cũng tò mò muốn xem cái `DeLorean' của em ấy nhanh đến đâu.''

Botan hất cằm tự hào, vỗ mạnh vào thành chiếc xe đẩy hàng của mình: ``Đừng có dám coi thường chiếc xe của ta nghe chưa!''

Nene đứng ngay sau lưng Botan, cau mày nhìn xuống chiếc xe mình đang đẩy: ``Xe nào của bà nữa? Đây là xe của tôi mà!''

Trong khi hai người ấy vẫn đang cãi nhau chí chóe về chiếc xe đẩy hàng nổi tiếng đó, Suisei và Marine đã chuyển sự chú ý sang một vấn đề quan trọng hơn nhiều. Ryuuhou khẽ nghiêng người, thì thầm vào tai Kanata: ``Em đặt cược 500 yên là Suisei-san sẽ hết hơi trước khi hết vòng đầu tiên, chị tham gia không?''

``Sao lại hỏi chị cái đó?'' nàng Thiên sứ toát mồ hôi hột. ``Chị cũng muốn kết thúc cái trò này lắm rồi\dots''

Nàng Ca sĩ hỏi cậu Tổng với vẻ mặt nghiêm chỉnh: ``Vậy chúng ta sẽ phân định thắng thua như thế nào ạ?''

Kuon nở một nụ cười đầy ẩn ý: ``Trong hai người, ai là người đầu tiên cán đích thì người đó thắng. Anh không cần hai em đua bằng cách nào, miễn là chạm được vạch đích trước là được.''

``Còn nếu\dots\ nếu người thua thì sao ạ?'' Marine hỏi lại một cách cẩn thận, giọng có chút run rẩy.

Kuon giơ lên ba ngón tay, tuyên bố: ``Người thua sẽ phải đứng ra đền bù toàn bộ chi phí sửa chữa văn phòng công ty. Con số đó cỡ khoảng\dots\ 30 triệu yên.''

``B-Ba mươi triệu yên!?'' cả hai người thốt lên kinh ngạc trước mức giá phi lý đó. Một làn gió lạnh lẽo bỗng thổi qua, khiến hai cô gái run lẩy bẩy như cầy sấy. Ryuuhou khẽ huýt sáo nhỏ, ghi nhanh con số 30 triệu lên tờ giấy của mình, thêm dấu chấm than to đậm để nhắc nhở.

``V-Vậy nếu\dots\ nếu Shishiron về đích trước thì sao ạ?'' Marine nuốt nước bọt ực một cái, hỏi bằng giọng đầy sợ hãi.

Kuon nở một nụ cười ``mất nhân tính'' : ``Thì hai người chia đôi số tiền đó, mỗi người đền 15 triệu yên.''

Nàng Sư tử bật cười lớn, vỗ mạnh vào vai cô bạn đồng hành, trong khi hai ``tay đua'' kia thì mặt xanh như tàu lá, y hệt vừa nhìn thấy hóa đơn tiền điện cuối tháng mà không có tiền trả. Ryuuhou gật gù như đồng tình với mức phạt, khẽ đạp nhẹ chân Kanata để cả hai cùng ngắm nhìn cảnh tượng hài hước trước mắt.

Không còn lối thoát nào khác, hai cô gái đành đạp xe chấp nhận tham gia cuộc đua định mệnh này.

Ba ``vận động viên'' đã đứng sẵn tại vạch xuất phát, trao cho nhau những ánh mắt đầy quyết tâm lửa cháy. Không khí căng thẳng bao trùm lấy toàn bộ khu vực, gương mặt ai cũng bừng bừng như sắp bước vào một trận chiến sinh tử.

Ryuuhou bước tới cạnh Kanata, trao cho bạn mình chiếc cờ xuất phát đã chuẩn bị sẵn, đồng thời nháy mắt khích lệ: ``Cứ phất mạnh thôi chị, nếu có ai chạy trước tín hiệu thì em phạt thua ngay.'' Kanata cầm lá cờ, bớt hồi hộp hơn một chút nhờ sự bình tĩnh của em gái.

``Chuẩn bị\dots'' Kanata hít một hơi thật sâu, giơ cao chiếc cờ lên trời theo lời nhắc của Ryuuhou.

``Sẵn sàng\dots'' Bầu không khí lập tức im lặng đến đáng sợ. Cả ba người đều dồn toàn bộ sự tập trung vào tín hiệu cuối cùng sắp vang lên, trong khi Ryuuhou đã bấm giờ sẵn trên điện thoại, sẵn sàng ghi nhận thời gian chính xác.

Và\dots

``\textbf{Bắt đầu!!}''

Kanata phất mạnh lá cờ xuống.

Ngay lập tức, bốn người Suisei, Marine, Botan và Nene cùng nhau lao về phía trước như ba cơn lốc. Đó không phải là cách nói cường điệu chút nào: sức mạnh của cú bứt phá quá khủng khiếp đến nỗi bụi bặm trên mặt đường bị thổi bay mù mịt, thậm chí Kanata đứng ở vạch xuất phát còn suýt bị cuốn theo dòng khí đó.

``Trời ơi\dots!'' cô nàng Thiên sứ vội vàng bám chặt vào cột đèn đường gần nhất để không bị ngã.

Kuon đứng kế bên, khoanh tay ngắm nhìn đám người vừa biến mất trên đường với vẻ mặt thích thú: ``Nhanh thật.''

Ryuuhou cười toe toét bên cạnh, tay vẫn cầm điện thoại ghi lại toàn bộ quá trình, trông như vừa thắng lớn một ván cờ. Cô vỗ nhẹ vào vai anh trai mình, giọng trêu chọc mà đầy tự tin: ``Em có bảo mà Onii. Mấy chị này lúc cãi nhau thì nói mình nhanh lắm, mà ra đường thì ai cũng không chịu kém ai. Em đặt cược 1000 yên là cả ba về đích cùng lúc, có dám nhận không?''

``Nhanh quá trời luôn!'' Kanata thở hổn hển, mắt vẫn không rời ba cái bóng đang nhỏ dần ở đường chân trời.

Nàng Trưởng nhóm nghiêng người sang, đưa màn hình điện thoại cho Kanata cùng xem, tốc độ của cả ba nhóm người hiện lên rõ rệt: ``Thôi nào Kanata-san, nhìn mà xem. Tốc độ của họ đang giữ gần như y hệt nhau. Suisei-san thì đạp chiếc tàu hỏa của mình không ngơi nghỉ, Marine-san cũng dốc sức với chiếc xe ba bánh, còn Botan-chan thì được Nene-san đẩy xe đẩy chạy cũng không kém. Làm sao mà có người về trước được?''

\ct

Cuộc đua trở nên căng thẳng tột độ. Đây không chỉ là cuộc chiến để xem ai là người nhanh nhất, mà còn là trận đấu sinh tử để tránh phải mất hàng chục triệu yên đền bù.

Suisei ghì chặt tay lái, răng nghiến chặt mà đạp hết sức: ``Không đời nào mình thua! Mình là `Tàu hỏa' mà, làm sao mà thua được!''

Marine cũng chẳng kém cạnh, vừa điều khiển chiếc xe ba bánh vừa thở dốc không ngừng: ``Chết thật, chân mình sắp rụng ra rồi! Hộc\dots\ Nhưng thà gãy chân còn hơn mất 30 triệu yên!!''

Trong khi đó, Botan đứng thẳng tắp trên chiếc xe đẩy hàng, hai tay dang rộng như đang tận hưởng cảm giác lướt gió, trông vô cùng thư thái: ``Ôi, cái cảm giác này đúng là `Fast \& Furry-ous' thật sự!''

Phía sau, Nene - người bị ép phải làm người đẩy xe cho cô - vừa chạy vừa thở hổn hển, mặt đỏ bừng như trái cà chua, mồ hôi nhễ nhại trên trán. Cô than van: ``B-Botan! Đừng có nằm đấy hưởng thụ nữa! Đạp phụ cho tôi đi chứ!''

``Không cần đâu, cứ tin vào `DeLorean' của chúng ta là được thôi,'' nàng Sư tử cười lớn, mái tóc bay phần phật trong gió.

``DeLorean cái gì mà DeLorean! Đây chỉ là xe đẩy hàng siêu thị thôi bà ơi!!''

Cả ba nhóm vẫn bám sát nhau không rời một ly. Đến đường Hamari, họ vẫn ngang ngửa. Vượt qua Đại lộ Mirete, vẫn không ai chịu nhường ai. Tất cả đều dốc hết sức lực, gió thổi rít qua tai, xen lẫn tiếng bánh xe lăn trên đường và tiếng chân chạy xào xạc.

Ryuuhou, người đã đạp xe đuổi kịp để ghi lại toàn bộ hành trình, đứng ở ngã tư Đại lộ, giơ điện thoại quay phim mà hét lớn khích lệ: ``Cố lên các chị ơi! Còn hai vòng nữa thôi! Ai về sau phải đền tiền đó nha!'' Cô nháy mắt trêu chọc, làm bầu không khí càng thêm sôi động.

Ba vòng đua cứ thế trôi qua, mỗi vòng cả ba đều dốc hết sức, không dám lơi lỏng chút nào. Ai cũng muốn cán đích đầu tiên, không chỉ để tránh bị phạt, mà còn để chứng minh mình mới là người nhanh nhất.

Chỉ vỏn vẹn một phút sau khi xuất phát, họ đã quay về lại điểm xuất phát.

``Ể? Đã xong rồi á!?'' Kanata trợn tròn mắt không thể tin được.

Ba chiếc xe lao qua vạch đích gần như cùng một lúc, bụi tung bay mù mịt đến mức không thể nhìn rõ người. Kuon phải lùi lại một bước, chớp mắt ngỡ ngàng trước khung cảnh hỗn loạn. Ryuuhou thì đứng ngay bên cạnh, ngừng quay video và bật cười lớn khi thấy cảnh tượng đó.

``Không thể nào\dots'' Cậu lẩm bẩm.

``Em có bảo mà Onii!'' Ryuuhou khoanh tay, đưa điện thoại cho Kuon và Kanata xem đoạn ghi hình vừa xong. ``1000 yên của em rồi đây, ai bảo không tin em.'' Cô nháy mắt tinh nghịch, rồi tiến tới kiểm tra tình hình của các chị.

Tất cả đều thở dốc, đổ gục xuống đất ngay khi xe dừng lại.

Kanata lại gần, cùng Ryuuhou xem lại đoạn video quay chậm trên điện thoại. Cô nhấp đi nhấp lại nhiều lần, mặt càng thêm ngạc nhiên.

``\dots Hòa rồi,'' cô nói, không giấu được sự bất ngờ. ``Cả ba người đều về đích cùng một lúc.''

Cô em gái của cậu Tổng gật đầu phụ họa, tay vẫn cầm điện thoại xem lại đoạn ghi hình: ``Em đoán đúng mà. Không ai chịu nhường ai được, khoảng cách giữa các xe chỉ có vài cm thôi. Làm sao mà phân định được ai thắng ai thua.''

``\textbf{Cái gì!??}'' Cả ba người bật dậy như bị lò xo đẩy lên.

``Làm sao mà hòa được!?'' Suisei hét lên.

``Khoan, vậy tức là không có ai thắng cả sao?'' Marine đổ mồ hôi lạnh.

``Đến tôi cũng không ngờ tới luôn\dots'' Nene vừa thở dốc vừa nói, mồ hôi vẫn nhiễu nhương trên người.

Kuon nhún vai, cười nhạt: ``Nếu hòa thì cả hai đứa đều thua. Nghĩa là\dots\ mỗi người đền 15 triệu yên.''

``\textbf{Không!!!!}''

Hai người lập tức ngã gục, mặt cắt không còn giọt máu. Họ không thể tin rằng sau bao nhiêu nỗ lực, lại không thể phân định được thắng thua. Điều đó có nghĩa là\dots\ họ vẫn phải chịu hình phạt!

Trong khi đó, Nene và Botan chỉ có thể cười trừ trước phản ứng quá mức của hai người kia. Nàng Sư tử vốn không bao giờ quan tâm đến thắng thua từ đầu, còn nàng Hải cẩu thì thở phào nhẹ nhõm khi nhận ra mình không phải chia sẻ số phận đó. Ryuuhou đứng cạnh họ cũng không nhịn được cười, nhưng cũng nhanh chóng lại gần để vỗ về hai người đàn chị.

``Thôi nào, hai chị đừng buồn nữa,'' cô nói, giọng vừa trêu chọc vừa có chút an ủi. ``Em biết là hai chị đã cố gắng hết sức rồi mà.''

``Thôi nào hai chị, đừng khóc lóc thế nữa,'' Botan cũng lại gần vỗ vai hai người đang gục xuống.

``\textbf{Nhưng mà đó là ba mươi triệu yên!!!}'' cả hai ngẩng mặt lên, mắt rưng rưng nhìn hậu bối.

Trước khi cơn tuyệt vọng nuốt chửng họ hoàn toàn, Nene bỗng nảy ra một ý: ``Khoan, vậy nếu chúng ta không trả bằng tiền thì sao?''

Ryuuhou cũng gật đầu đồng tình, mắt lóe lên sáng suốt: ``Đúng đó! Chúng ta có thể kêu gọi mọi người trong công ty cùng nhau xây lại văn phòng mà. Chẳng phải trước đây cũng từng làm như vậy rồi sao? Mọi người cùng góp sức, chẳng ai phải bỏ ra nhiều tiền cả.''

Từ xa, Kuon nhún vai đáp: ``Thì đền bằng hiện vật. Cùng nhau tìm cách xây lại cả tòa nhà đó.''

Một tia hy vọng bùng lên trong mắt hai cô gái, như thể vừa tìm thấy ánh sáng cuối đường hầm tối tăm.

``Vậy là chỉ cần kêu gọi mọi người cùng xây lại là xong đúng không!?'' Marine nhanh chóng nắm bắt được vấn đề.

``Vậy thì dễ quá!'' Suisei cũng hào hứng hẳn lên, khẽ thở phào vì cuối cùng cũng không phải nộp ra số tiền khổng lồ kia.

\ct

Chẳng mấy chốc, cả đoàn đã kéo toàn bộ nhân viên và các tài năng trong công ty đến chung tay hỗ trợ. Một chiến dịch tổng lực được phát động, biến tầng văn phòng tan hoang sau vụ nổ thành một công trường xây dựng nhộn nhịp.

Aqua được giao nhiệm vụ dọn dẹp đống gạch vụn chồng chất. Subaru đảm nhận vai trò trộn xi măng. Còn Noel thì tự nhiên biến thành người chuyên vác những tấm gỗ khổng lồ, sức mạnh phi thường đến mức chẳng cần dùng đến bất kỳ dụng cụ hỗ trợ nào. Riêng với Pekora\dots\ chẳng ai hiểu nổi tại sao cô nàng lại phải đóng vai thợ hồ, trong khi rõ ràng vụ nổ vừa rồi chẳng hề dính líu gì đến cô cả.

``Khoan đã, tại sao tôi lại phải làm cái việc này, peko!?'' cô nàng không khỏi hét lên bất bình, tay vẫn đang cầm xẻng xúc xi măng bực bội.

``Vì em là người hay gây ra chuyện lộn xộn nhất cả công ty mà,'' Kuon vừa bụm miệng cười trộm vừa đáp lời.

``E-em có làm gì đâu chứ, peko!!'' Pekora vẫn cố gắng phân trần.

``Làm đi nào, đừng nói nữa,'' Kanata vừa nói vừa búng trán cô một cái rõ đau.

Pekora chỉ biết ôm trán rên rỉ, uất ức trong lòng nhưng chẳng dám cãi lại ai. Cuối cùng, nhờ vào sự cố gắng không ngừng nghỉ của tất cả mọi người - kể cả cô thợ hồ Pekora bất đắc dĩ - toàn bộ văn phòng đã được xây dựng lại nguyên vẹn, mà không ai phải tốn một đồng xu nào ra cả.

Khi mọi người cùng nhau nghỉ ngơi sau buổi lao động vất vả, Ryuuhou ngồi xuống cạnh anh trai mình, giọng đầy tinh nghịch: ``Onii, em phải thú thật là vụ nổ lần này đâu phải do em gây ra đâu. Lúc đó em đang ở trên trời cùng với Suisei-san, Marine-san và Kanata-san mà. Em chỉ kịp thấy văn phòng bốc khói, rồi cả đám cứ thế bay theo cơn nổ lên trời luôn.''

Kuon nhìn em gái, nhíu mày hỏi: ``Thế em có bị sao không?''

``Không có gì đâu ạ, vì lúc đó em đang ngồi trên chiếc `tàu hỏa' của Suisei nên được bay theo lên trời mất rồi. Cũng may mắn lắm, chứ nếu em ở dưới đất thì giờ này chắc cũng phải đứng cùng hàng với mấy chị kia để chia tiền đền bù mất.''

Kuon chỉ biết cười xoa đầu em gái: ``Biết thế thì lần sau thấy có dấu hiệu gì sắp nổ thì nhảy lên tàu của Suisei ngay nhé.''

``Em nhớ rồi. Mà Onii này, em thấy trò Old Maid nguy hiểm thật đấy. Lần sau chúng ta chơi bài thì đổi trò khác đi nhé.''

``Ờ, đồng ý.''

Tuy nhiên, vẫn có một điều chẳng bao giờ thay đổi\dots

``Dù có nói gì đi nữa thì người nhanh nhất vẫn là tôi!'' Suisei không chịu thua, tiếp tục tranh cãi.

``Điều gì nhảm nhí vậy! Tôi cán đích trước cô cả nửa cái bánh xe đấy nhé!'' Marine lập tức phản bác lại.

``Đừng có quên là chúng ta còn có `DeLorean', hai cô còn thua xa tôi đấy!'' Botan cũng cười lớn chen vào.

Và thế là, cuộc tranh luận vô hồi kết lại tiếp tục diễn ra, chẳng ai chịu nhường ai. Chỉ có một điểm duy nhất mà cả ba người này đều đồng tình với nhau\dots

Họ cùng thề sẽ không bao giờ đụng đến trò Old Maid một lần nữa.

``Đấy đúng là yêu cầu hợp lý của anh đấy, Kuon-senpai,'' Kanata thở dài thườn thượt.

Kuon cũng gật gù đồng tình: ``Ừm, trò chơi ấy dễ làm mọi người mất đoàn kết lắm. Nhiều lúc còn gây ra thương tích nữa.''

Ryuuhou cũng thêm vào với giọng đùa cợt: ``Em còn nghe nói hồi trước có vụ Noel-san cắm cái chùy vào mông Subaru-san nữa. Không biết lúc đó có ai quay lại cảnh đó không, em xin làm tư liệu xem sao.''

Cả nhóm lập tức im lặng, ai cũng nhớ lại khoảnh khắc kinh hoàng đó, đồng loạt rùng mình: ``Ừ thôi, đừng nhắc lại chuyện đó nữa\dots'' / ``Ờm, em không cần phải biết đâu.''

Ryuuhou chỉ biết cười khúc khích, sau đó đứng dậy phủi bụi trên quần: ``Em đi pha trà đây. Ai uống không? Uống xong rồi còn phải đi dọn đống bìa các-tông của Suisei-san nữa.''

``Cho anh một cốc,'' Kuon lập tức nói.

``Chị cũng thế,'' Kanata cũng giơ tay đăng ký.

Và thế là, một ngày nữa lại trôi qua trong chính sự hỗn loạn thường nhật của \hll. Mùa hè đang đến gần từng ngày, và Kuon thừa biết những ngày sắp tới sẽ còn vô vàn chuyện bất ngờ không thể đoán trước nữa. Nhưng có một điều anh luôn chắc chắn: dù có chuyện gì xảy ra, cậu và Ryuuhou - cùng với tất cả những người bạn của mình - sẽ vẫn luôn ở bên cạnh nhau.

Ngoài ô cửa sổ, ánh nắng cuối chiều đang nhuộm vàng cả một vùng bầu trời. Những đám mây trắng trôi lững lờ trên cao, và một làn gió nhẹ mang theo hơi ấm của mùa hè đang đến rất gần.

\end{document}